\begin{center}
	\zihao{4}\heiti
	% 中国科学院大学\\
	\LARGE{中国科学院研究生院\\$\boldsymbol{2008}$年招收硕士研究生入学统一考试试卷\\
	科目名称：电动力学}

\end{center}
\noindent
考生须知：\\
1.本试卷满分为150分,全部考试时间总计180分钟。\\
2.所有答案必须写在答题纸上,写在试题纸上或草稿纸上一律无效。

\vspace*{2ex}
\hrule
\vspace*{4ex}

\fancypagestyle{mypagestyle}{%
	\fancyhf{}% Clear header/footer
	\fancyfoot[OR]{第~\thepage~页，共2页}%
	\fancyfoot[EL]{第~\thepage~页，共2页}%
	\fancyfoot[OL]{科目名称：电动力学 }%
	\fancyfoot[ER]{科目名称：电动力学 }%
	\renewcommand{\headrulewidth}{0pt}% Header rule of .4pt
	\renewcommand{\footrulewidth}{0.5pt}% Header rule of .4pt
}
\fancypagestyle{plain}{%
	\fancyhf{}% Clear header/footer
	\fancyfoot[OR]{第~\thepage~页，共2页}%
	\fancyfoot[EL]{第~\thepage~页，共2页}%
	\fancyfoot[OL]{科目名称：电动力学 }%
	\fancyfoot[ER]{科目名称：电动力学 }%
}

\setcounter{page}{1}%设置页面计数器
\pagestyle{mypagestyle}

\onehalfspacing
% 排版结束，下面是题目
\vspace{0.4cm}
\noindent
一、问答题
\vspace{0.2cm}


\noindent
1.\hspace{0.8em}写出$\rho \neq 0$，$J \neq 0$情况下电场强度和磁场强度满足的波动方程
\vspace{0.1cm}

\noindent
2.\hspace{0.8em}写出在$\Vec{E}$垂直于入射面时，电磁波在全反射时电场的反射系数和相移公式
\vspace{0.1cm}

\noindent
3.\hspace{0.8em}写出球坐标系下电偶极辐射的远区电磁场分布
\vspace{0.1cm}

\noindent
4.\hspace{0.8em}写出运动电荷的李纳-维谢尔势

\vspace{-1.6cm}
\begin{figure}[h!]
    \flushright
    \includegraphics[height=5cm]{test/2008/pic_1.png}
\end{figure}

\vspace{-0.8cm}
\noindent
二、\hspace{0.2em}选择题
\vspace{0.2cm}

\noindent
1.\hspace{0.8em}在介电系数为$\boldsymbol{\varepsilon}$，无限大各向同性均匀介质中，挖出一个半径为$R_0$的球形真空腔，并把整个系统放在一个均匀的外电场$\Vec{E_0}$中，则空腔外一点A和空腔内一点B的电场强度满足
\vspace{0.2cm}

A.\hspace{0.8em}$E_A = 0, E_B \neq 0$  \hspace{1em} B.\hspace{0.8em}$E_A \neq 0, E_B = 0$ \hspace{1em} C.\hspace{0.8em}$E_A < E_B $ \hspace{1em} D.\hspace{0.8em}$E_A > E_B $ 
\vspace{0.2cm}

\noindent
2.\hspace{0.8em}均匀带电的旋转椭圆球体带总电量为Q，坐标原点在椭球中心处，其
\vspace{0.2cm}

A.\hspace{0.8em}电偶极矩和电四极矩与$Q^2$有关 \hspace{2em} B.\hspace{0.8em}电偶极矩为零，电四极矩不为零\vspace{0.1cm}

C.\hspace{0.8em}电偶极矩不为零，电四极矩为零 \hspace{2em} D.\hspace{0.8em}电偶极矩，电四极矩都为零

\vspace{0.2cm}

\noindent
3.\hspace{0.8em}某粒子的总能量是它的静止能量的$k$倍，则其运动速度的大小为
\vspace{0.3cm}

A.\hspace{0.8em}$\displaystyle{ \frac{c}{k-1} }$  \hspace{2em} B.\hspace{0.8em}$\displaystyle{ \frac{c}{k} \sqrt{1-k^2} }$   \hspace{2em} C.\hspace{0.8em}$\displaystyle{ \frac{c}{k} \sqrt{k^2-1} }$   \hspace{2em} D.\hspace{0.8em}$\displaystyle{ \frac{c}{k+1} \sqrt{k(k+2)} }$

\vspace{0.3cm}

\noindent
4.\hspace{0.8em}当两种导电介质内流有稳恒电流时，分界面上的电力线与界面法线夹角满足
\vspace{0.3cm}

A.\hspace{0.8em}$\displaystyle{\frac{\tan \theta_1}{\tan \theta_1} = \frac{\boldsymbol{\varepsilon_2}}{\boldsymbol{\varepsilon_1}}}$ \hspace{2em} B.\hspace{0.8em}$\displaystyle{\frac{\tan \theta_1}{\tan \theta_1} = \frac{\boldsymbol{\varepsilon_1}}{\boldsymbol{\varepsilon_2}}}$ \hspace{2em} C.\hspace{0.8em}$\displaystyle{\frac{\tan \theta_1}{\tan \theta_1} = \frac{\boldsymbol{\sigma_2}}{\boldsymbol{\sigma_1}}}$ \hspace{2em} D.\hspace{0.8em}$\displaystyle{\frac{\tan \theta_1}{\tan \theta_1} = \frac{\boldsymbol{\sigma_1}}{\boldsymbol{\sigma_2}}}$ 
\vspace{0.3cm}

\noindent
5.\hspace{0.8em}$\sqrt{\boldsymbol{\varepsilon_0} / \mu_0}$物理计量单位为
\vspace{0.2cm}

A.\hspace{0.8em} 欧\hspace{2em}
B.\hspace{0.8em} 西\hspace{2em}
C.\hspace{0.8em} 法\hspace{2em}
D.\hspace{0.8em} 享
\vspace{0.2cm}

\noindent
6.\hspace{0.8em}柱坐标下磁化强度为$\Vec{M} = m \cos \theta \, \Vec{e_r}$，磁化电流密度为
\vspace{0.3cm}

A.\hspace{0.8em} $\displaystyle{ \frac{m}{r} \sin \theta \, \Vec{e_z} }$\hspace{2em}
B.\hspace{0.8em} $ m \sin \theta \,  \Vec{e_z} $\hspace{2em}
C.\hspace{0.8em} $\displaystyle{ \frac{m}{r} \sin \theta \, \Vec{e_{\theta}} }$\hspace{2em}
D.\hspace{0.8em} $\displaystyle{ \frac{m}{r} \cos \theta \, \Vec{e_z} }$
\vspace{0.3cm}

\noindent
7.\hspace{0.8em}电磁波在良导体内衰减传播，其磁场相位比电场相位超前的角度为
\vspace{0.2cm}

A.\hspace{0.8em} $-90^{\circ}$ \hspace{2em}
B.\hspace{0.8em} $90^{\circ}$\hspace{2em}
C.\hspace{0.8em} $45^{\circ}$\hspace{2em}
D.\hspace{0.8em} $-45^{\circ}$

\newpage
\noindent
8.\hspace{0.8em}D是四电极距，下面几个表达式中哪个具有功率的量纲？（此题不一定存在正确答案
）
\vspace{0.3cm}

A.\hspace{0.4em} $\displaystyle{ \frac{\mu_0}{4 \pi} \frac{|\ddot{D}|^2}{540C^3} }$\hspace{1em}
B.\hspace{0.4em} $\displaystyle{ \frac{1}{4 \pi \boldsymbol{\varepsilon_0} } \frac{|\ddot{D}|^2}{540C^4} }$\hspace{1em}
C.\hspace{0.4em} $\displaystyle{ \frac{1}{4 \pi \boldsymbol{\varepsilon_0} } \frac{|\ddot{D}|^2}{540C^5} }$\hspace{1em}
D.\hspace{0.4em} $\displaystyle{ \frac{1}{4 \pi \boldsymbol{\varepsilon_0} } \frac{|\ddot{D}|}{540C^4} }$\hspace{1em}
E.\hspace{0.4em} $\displaystyle{ \frac{1}{4 \pi \boldsymbol{\varepsilon_0} } \frac{|\ddot{D}|^2 \omega^2}{540C^4} }$
\vspace{0.5cm}

\noindent
三、\hspace{0.2em}有一均匀带电的导体薄球壳，其总电量为Q,半径为$R_0$，当球壳绕自身某一直径以角速度$\omega$转动时，试用磁标势求球内外的磁感应强度B

\vspace{5cm}

\noindent
四、\hspace{0.2em}有一沿z轴无限长波导管，其截面为矩形，边长为a和b，腔内介电常数为$\varepsilon$，磁导率为$\mu$，求其最低频率的横磁模式(TM)电磁波的电磁场强度$\Vec{E}$和$\Vec{H}$各分量，以及最低震荡频率$\omega_0$

\vspace{5cm}

\noindent
五、\hspace{0.2em}如图所示，设折射率为n的均匀透明流体介质中有一固定的点光源S，
流体介质速度$\Vec{v}$均匀地沿SP方向流动，其速度的大小与该介质内的光速相同，在与SP成$\theta$角的方向，距离为$l_0$外放有一个很小的光线接收器R。假设S和R都对流体的流动无阻碍，考虑相对论效应，求光线到达接收器所需要的时间，并与非相对论的结果进行比较(图中$\Vec{v}$为可以到达接收器的光在静止介质内的传播速度)
\vspace{1cm}
\begin{figure}[h!]
    \flushright
    \includegraphics[height=6cm]{test/2008/pic_2.png}
\end{figure}